% Options for packages loaded elsewhere
\PassOptionsToPackage{unicode}{hyperref}
\PassOptionsToPackage{hyphens}{url}
\documentclass[
  ignorenonframetext,
]{beamer}
\newif\ifbibliography
\usepackage{pgfpages}
\setbeamertemplate{caption}[numbered]
\setbeamertemplate{caption label separator}{: }
\setbeamercolor{caption name}{fg=normal text.fg}
\beamertemplatenavigationsymbolsempty
% remove section numbering
\setbeamertemplate{part page}{
  \centering
  \begin{beamercolorbox}[sep=16pt,center]{part title}
    \usebeamerfont{part title}\insertpart\par
  \end{beamercolorbox}
}
\setbeamertemplate{section page}{
  \centering
  \begin{beamercolorbox}[sep=12pt,center]{section title}
    \usebeamerfont{section title}\insertsection\par
  \end{beamercolorbox}
}
\setbeamertemplate{subsection page}{
  \centering
  \begin{beamercolorbox}[sep=8pt,center]{subsection title}
    \usebeamerfont{subsection title}\insertsubsection\par
  \end{beamercolorbox}
}
% Prevent slide breaks in the middle of a paragraph
\widowpenalties 1 10000
\raggedbottom
\AtBeginPart{
  \frame{\partpage}
}
\AtBeginSection{
  \ifbibliography
  \else
    \frame{\sectionpage}
  \fi
}
\AtBeginSubsection{
  \frame{\subsectionpage}
}
\usepackage{iftex}
\ifPDFTeX
  \usepackage[T1]{fontenc}
  \usepackage[utf8]{inputenc}
  \usepackage{textcomp} % provide euro and other symbols
\else % if luatex or xetex
  \usepackage{unicode-math} % this also loads fontspec
  \defaultfontfeatures{Scale=MatchLowercase}
  \defaultfontfeatures[\rmfamily]{Ligatures=TeX,Scale=1}
\fi
\usepackage{lmodern}
\ifPDFTeX\else
  % xetex/luatex font selection
\fi
% Use upquote if available, for straight quotes in verbatim environments
\IfFileExists{upquote.sty}{\usepackage{upquote}}{}
\IfFileExists{microtype.sty}{% use microtype if available
  \usepackage[]{microtype}
  \UseMicrotypeSet[protrusion]{basicmath} % disable protrusion for tt fonts
}{}
\makeatletter
\@ifundefined{KOMAClassName}{% if non-KOMA class
  \IfFileExists{parskip.sty}{%
    \usepackage{parskip}
  }{% else
    \setlength{\parindent}{0pt}
    \setlength{\parskip}{6pt plus 2pt minus 1pt}}
}{% if KOMA class
  \KOMAoptions{parskip=half}}
\makeatother
\setlength{\emergencystretch}{3em} % prevent overfull lines
\providecommand{\tightlist}{%
  \setlength{\itemsep}{0pt}\setlength{\parskip}{0pt}}
% Auto-generated by infrastructure/rendering/_pdf_latex_helpers.py::extract_math_font_preamble.
% Minimal header passed to Pandoc via -H so Beamer slides pick up the same math font
% as the combined-PDF path. Adjust the manuscript's preamble.md to change the font globally.
\usepackage{unicode-math}
\setmathfont{latinmodern-math.otf}

% Slide layout defaults for warning-clean scientific decks.
\usepackage{etoolbox}
\IfFileExists{xurl.sty}{\usepackage{xurl}}{}
\IfFileExists{seqsplit.sty}{\usepackage{seqsplit}}{\newcommand{\seqsplit}[1]{#1}}
\protected\def\breakseq#1{\seqsplit{#1}}
\protected\def\breaktt#1{\begingroup\ttfamily\seqsplit{#1}\endgroup}
\setlength{\emergencystretch}{6em}
\tolerance=5000
\hbadness=10000
\hfuzz=1pt
\setlength{\tabcolsep}{2pt}
\AtBeginEnvironment{longtable}{\tiny\renewcommand{\arraystretch}{0.86}\setlength{\tabcolsep}{1pt}}
\AtBeginEnvironment{tabular}{\tiny\renewcommand{\arraystretch}{0.86}\setlength{\tabcolsep}{1pt}}

% Natbib and cross-reference fallbacks — slides don't load natbib
% or cleveref, but combined-PDF manuscript prose may emit these
% commands. The fallback renders citations as a bracketed key list
% and cross-references as detokenized labels so slides don't fail on
% undefined control sequences or raw underscores. \providecommand is
% a no-op if packages load later.
\providecommand{\citep}[1]{[#1]}
\providecommand{\citet}[1]{#1}
\providecommand{\citealp}[1]{#1}
\providecommand{\citeauthor}[1]{#1}
\providecommand{\citeyear}[1]{#1}
\providecommand{\cref}[1]{\texttt{\detokenize{#1}}}
\providecommand{\Cref}[1]{\texttt{\detokenize{#1}}}
\usepackage{bookmark}
\IfFileExists{xurl.sty}{\usepackage{xurl}}{} % add URL line breaks if available
\urlstyle{same}
% fallback for those not using the hyperref driver hyperxmp:
\makeatletter
\@ifundefined{xmpquote}{\newcommand{\xmpquote}[1]{#1}}{}
\makeatother
\hypersetup{
  hidelinks,
  pdfcreator={LaTeX via pandoc}}

\author{\texorpdfstring{}{}}
\date{}

\begin{document}

\section{Scope, Related Work, and Positioning}\label{sec:scope}

\begin{frame}[allowframebreaks]{Scope, Related Work, and Positioning}
This section situates the exemplar and states explicit boundaries. The
goal is not to compete with BPL's full compiler pipeline \breakseq{[@bpl2026]}
--- a \textasciitilde32,000-line implementation with a Lark grammar, a
robot backend, and a hash-chained audit/compliance layer --- but to show
how a minimal, test-backed subset of BPL's domain-language design fits
the template's reproducibility and rendering stack
\breakseq{[@peng2011reproducible]}, generalized from wet-lab protocols to any
controlled procedure.
\end{frame}

\begin{frame}[allowframebreaks]{Domain-specific languages for controlled
procedures}
\protect\phantomsection\label{domain-specific-languages-for-controlled-procedures}
Encoding a procedure as a program rather than free text is a
long-standing software-engineering pattern: Fowler's treatment of
domain-specific languages \breakseq{[@fowler2010dsl]} frames the general case
--- a notation restricted to exactly the concepts a domain needs. BPL
\breakseq{[@bpl2026]} applies this specifically to laboratory protocols, adding
biology-native types (reagents, labware, MW-aware concentrations),
staged validation, and deterministic compilation to a robot or human
execution target. The present manuscript restricts attention to the
parts of that design that generalize beyond biology: a controlled
vocabulary of step intents, a small dimensional-safety unit system,
staged validation gates, and deterministic compilation to a hashed plan.
\end{frame}

\begin{frame}[fragile,allowframebreaks]{What this exemplar does not
implement (out of scope)}
\protect\phantomsection\label{what-this-exemplar-does-not-implement-out-of-scope}
\begin{enumerate}
\tightlist
\item
  \textbf{A text grammar and parser.} BPL parses \texttt{.bpl} source
  through a Lark grammar into a typed AST. This exemplar constructs a
  \texttt{Method} directly as frozen Python dataclasses --- no concrete
  syntax, no parser, no AST layer.
\item
  \textbf{Biology-native types.} BPL's unit system includes MW-aware
  concentration conversions and reagent physical-form metadata
  (\texttt{cas}, \texttt{physical\_form}). This exemplar's
  \texttt{units.py} implements only the dimensional-safety subset (mass,
  volume, temperature, time, concentration, count, dimensionless) needed
  to demonstrate the ``the type system catches \texttt{mL\ +\ g}''
  guarantee.
\item
  \textbf{A robot backend.} BPL lowers intents to Biomek i7 primitives
  (\texttt{aspirate}, \texttt{dispense}, \texttt{pick\_tips}). This
  exemplar's \breaktt{Target.AUTOMATED} has no backend-specific lowering
  stage; it is a scheduling and gate-compatibility concept only.
\item
  \textbf{Cryptographic audit guarantees.} \texttt{trust.py}'s
  hash-chain is deliberately scoped to the same honest boundary BPL
  itself claims: a consistency check against accidental corruption, not
  a tamper-proof guarantee against an actor with write access to the
  entire chain.
\item
  \textbf{SOP-to-DSL generation.} BPL includes an agentic workflow that
  translates natural-language SOPs into validated \texttt{.bpl}
  programs. This exemplar's worked examples are hand-authored Python,
  not generated.
\end{enumerate}
\end{frame}

\begin{frame}[fragile,allowframebreaks]{What this project proves about
the template}
\protect\phantomsection\label{what-this-project-proves-about-the-template}
The validation and compilation steps here are a deliberately small
subset of BPL's. The \textbf{non-standard} contribution is procedural:
the same tested functions in \breaktt{src/methods\_dsl/} back the
analysis script, the test suite, and this manuscript, so the
compiled-plan table and the figure always refer to the same code. That
pattern --- and the specific generalization from a biology-only domain
language to a domain-neutral one --- is what downstream projects should
copy, whether the controlled procedure is a wet-lab protocol, an
instrument calibration sweep, or a computational pipeline.
\end{frame}

\begin{frame}[fragile,allowframebreaks]{Explicit limitations}
\protect\phantomsection\label{explicit-limitations}
\begin{enumerate}
\tightlist
\item
  \textbf{Two worked examples}: \texttt{PBSPreparation} and
  \breaktt{SensorCalibrationSweep} exercise every gate-success path but
  are not a corpus; gate-failure coverage instead lives in
  \breaktt{tests/conftest.py}'s dedicated fixtures.
\item
  \textbf{No backend lowering}: \texttt{Target} selects a compatibility
  class, not a concrete execution backend; no robot or simulation
  runtime exists in this exemplar.
\item
  \textbf{Unit table, not a full unit library}: seven dimensions and a
  fixed unit table, not a general-purpose dimensional-analysis library
  like \texttt{pint}.
\item
  \textbf{No persistence layer}: \texttt{trust.py}'s hash-chain lives in
  memory for the duration of one script run; this exemplar does not
  implement durable storage for a real audit trail.
\end{enumerate}

These limitations are intentional: they narrow the surface so that the
reproducibility concerns --- tested functions, a thin script, and
generated-variable-injected prose --- remain visible rather than buried
under a compiler implementation at BPL's full scale.
\end{frame}

\end{document}
